\documentclass[a4paper]{article}
\usepackage[fontset=fandol]{ctex}
\usepackage{amsmath, amssymb}
\usepackage{graphicx}
\usepackage[margin=2.5cm]{geometry}
\usepackage[most]{tcolorbox}
\usepackage{etoolbox}
\usepackage{listings}
\usepackage{hyperref}
\usepackage{booktabs}
\usepackage{subcaption}
\usepackage{float}
\usepackage{tikz}
\IfFileExists{pgfplots.sty}{\usepackage{pgfplots}\pgfplotsset{compat=1.18}}{}
\newtcolorbox{knowledgebox}[1]{enhanced, colback=blue!5!white, colframe=blue!75!black, colbacktitle=blue!75!black, coltitle=white, fonttitle=\bfseries, title=#1, attach boxed title to top left={yshift=-2mm, xshift=2mm}, boxrule=1pt, sharp corners}
\newtcolorbox{importantbox}[1]{enhanced, colback=yellow!10!white, colframe=yellow!80!black, colbacktitle=yellow!80!black, coltitle=black, fonttitle=\bfseries, title=#1, sharp corners}
\newtcolorbox{warningbox}[1]{enhanced, colback=red!5!white, colframe=red!75!black, colbacktitle=red!75!black, coltitle=white, fonttitle=\bfseries, title=#1, sharp corners}
\lstset{language=Python,basicstyle=\ttfamily\small,keywordstyle=\color{blue},stringstyle=\color{red!60!black},commentstyle=\color{green!60!black},breaklines=true,frame=single,numbers=left,numberstyle=\tiny\color{gray},captionpos=b,extendedchars=false}
\newcommand{\notetitle}{Acoustic Model Overview}
\newcommand{\noteauthors}{Codex 基于课程页、公开视频字幕、coverage ledger 与补充官方材料整理}
\newcommand{\notedate}{\today}
\newcommand{\videochannel}{WAVLab}
\newcommand{\videopublishdate}{2023-09-13}
\newcommand{\videoduration}{unknown}
\newcommand{\videourl}{https://www.youtube.com/watch?v=GTaqSmQSHBs}
\newcommand{\videocoverpath}{cover.jpg}
\begin{document}
\begin{titlepage}
\centering
{\Large 课程笔记\par}
\vspace{1.2cm}
{\huge\bfseries \notetitle\par}
\vspace{0.8cm}
{\large \noteauthors\par}
\vspace{0.3cm}
{\large \notedate\par}
\vspace{1.2cm}
\ifdefempty{\videocoverpath}{{\small 公开视频封面不可用。\par}}{\includegraphics[width=0.82\textwidth,height=0.45\textheight,keepaspectratio]{\videocoverpath}\par}
\vfill
\begin{tcolorbox}[width=0.9\textwidth, colback=black!2!white, colframe=black!60, sharp corners]
\textbf{视频作者/频道}：\videochannel\par
\textbf{发布日期}：\videopublishdate\par
\textbf{视频时长}：\videoduration\par
\textbf{视频链接}：\href{\videourl}{\nolinkurl{\videourl}}
\end{tcolorbox}
\end{titlepage}
\tableofcontents
\newpage
\begin{warningbox}{Source Discrepancy}
2023 course page lists Acoustic Model Overview but no public 2023 recording was found. This chapter is reconstructed from the 2023 schedule row plus official WAVLab Fall 2022 same-course materials on advanced acoustic modeling, deep acoustic models, and system overview.
\end{warningbox}

\section{5.0 重建前言：为什么需要 Acoustic Model Overview}
根据课程页主题、公开视频字幕以及本讲的 coverage ledger，5.0 重建前言：为什么需要 Acoustic Model Overview 这一部分主要围绕 2023 course row: Acoustic Model Overview | role of the acoustic model in an ASR system 展开。
写作上保持 coverage-first 原则：先把老师在这一阶段真正讲过的问题设定、模型部件、训练或推理逻辑讲清，再压缩成适合复习的结构。

\begin{figure}[H]
\centering
\includegraphics[width=0.92\textwidth]{figures/coverage_map.png}
\caption{本讲 coverage map：按 coverage\_units.jsonl 归纳出的核心主题与章节映射。}
\end{figure}

\begin{figure}[H]
\centering
\includegraphics[width=0.92\textwidth]{figures/timeline.png}
\caption{本讲 timeline：根据字幕时间范围归纳的主题推进顺序。}
\end{figure}

\subsection{2023 course row: Acoustic Model Overview}
从证据位置看，这一部分对应的课程证据区间是 2023-09-13。课程在这里强调的核心内容可以概括为：2023 course row: Acoustic Model Overview | role of the acoustic model in an ASR system。
从课程组织上看，这一部分承担的是把局部概念嵌回整条识别流水线中的作用，避免读者把公式或模块孤立地记忆。

记录系统中的附加说明指出：This section is reconstructed from the 2023 course row and official same-course auxiliary materials because no public 2023 recording or slides were found.

在本章结构中，这一单元被并入“5.0 重建前言：为什么需要 Acoustic Model Overview”，目的是把同一阶段的定义、机制和权衡集中讲清。

\subsection{本章小结}
5.0 重建前言：为什么需要 Acoustic Model Overview 这一部分把 Acoustic Model Overview 中相邻的教学单元串成了一条连续叙事：既交代了这一阶段要解决的问题，也交代了方法、推理或工程权衡为什么要这样组织。

\section{5.1 系统视角下的 Acoustic Model}
根据课程页主题、公开视频字幕以及本讲的 coverage ledger，5.1 系统视角下的 Acoustic Model 这一部分主要围绕 从系统视角比较 HMM-based systems 与 end-to-end systems，并解释 acoustic model 在整条识别流水线中的位置。 展开。
写作上保持 coverage-first 原则：先把老师在这一阶段真正讲过的问题设定、模型部件、训练或推理逻辑讲清，再压缩成适合复习的结构。

\subsection{从系统视角比较 HMM-based systems 与 end-to-end systems，并解释 acoustic model 在整条识别流水线中的位置}
从证据位置看，这一部分对应的课程证据区间是 00:00:00,000--00:12:00,000。课程在这里强调的核心内容可以概括为：从系统视角比较 HMM-based systems 与 end-to-end systems，并解释 acoustic model 在整条识别流水线中的位置。。
这里的重点不是把 HMM 当作过时历史，而是把它当作对齐、状态建模与动态规划最清晰的教学载体。

记录系统中的附加说明指出：Grounded in official Fall 2022 same-course materials; used as auxiliary reconstruction for the missing 2023 public session.

在本章结构中，这一单元被并入“5.1 系统视角下的 Acoustic Model”，目的是把同一阶段的定义、机制和权衡集中讲清。

\subsection{本章小结}
5.1 系统视角下的 Acoustic Model 这一部分把 Acoustic Model Overview 中相邻的教学单元串成了一条连续叙事：既交代了这一阶段要解决的问题，也交代了方法、推理或工程权衡为什么要这样组织。

\section{5.2 经典 Advanced Acoustic Modeling}
根据课程页主题、公开视频字幕以及本讲的 coverage ledger，5.2 经典 Advanced Acoustic Modeling 这一部分主要围绕 phonetic decision tree、context-dependent states 与 adaptation 为何是 advanced acoustic modeling 的核心。 展开。
写作上保持 coverage-first 原则：先把老师在这一阶段真正讲过的问题设定、模型部件、训练或推理逻辑讲清，再压缩成适合复习的结构。

\subsection{phonetic decision tree、context-dependent states 与 adaptation 为何是 advanced acoustic modeling 的核心}
从证据位置看，这一部分对应的课程证据区间是 00:12:00,000--00:25:00,000。课程在这里强调的核心内容可以概括为：phonetic decision tree、context-dependent states 与 adaptation 为何是 advanced acoustic modeling 的核心。。
从课程组织上看，这一部分承担的是把局部概念嵌回整条识别流水线中的作用，避免读者把公式或模块孤立地记忆。

记录系统中的附加说明指出：Auxiliary reconstruction using official Fall 2022 Advanced Acoustic Modeling video and same-course page row.

在本章结构中，这一单元被并入“5.2 经典 Advanced Acoustic Modeling”，目的是把同一阶段的定义、机制和权衡集中讲清。

\subsection{本章小结}
5.2 经典 Advanced Acoustic Modeling 这一部分把 Acoustic Model Overview 中相邻的教学单元串成了一条连续叙事：既交代了这一阶段要解决的问题，也交代了方法、推理或工程权衡为什么要这样组织。

\section{5.3 从 GMM-HMM 到 DNN Acoustic Model}
根据课程页主题、公开视频字幕以及本讲的 coverage ledger，5.3 从 GMM-HMM 到 DNN Acoustic Model 这一部分主要围绕 从 GMM-HMM 过渡到 DNN acoustic model，解释 posterior、state targets 与 hybrid recipe。 展开。
写作上保持 coverage-first 原则：先把老师在这一阶段真正讲过的问题设定、模型部件、训练或推理逻辑讲清，再压缩成适合复习的结构。

\subsection{从 GMM-HMM 过渡到 DNN acoustic model，解释 posterior、state targets 与 hybrid recipe}
从证据位置看，这一部分对应的课程证据区间是 00:25:00,000--00:40:00,000。课程在这里强调的核心内容可以概括为：从 GMM-HMM 过渡到 DNN acoustic model，解释 posterior、state targets 与 hybrid recipe。。
这里的重点不是把 HMM 当作过时历史，而是把它当作对齐、状态建模与动态规划最清晰的教学载体。

记录系统中的附加说明指出：Auxiliary reconstruction from official same-course materials on advanced acoustic modeling and deep learning for speech recognition.

在本章结构中，这一单元被并入“5.3 从 GMM-HMM 到 DNN Acoustic Model”，目的是把同一阶段的定义、机制和权衡集中讲清。

\subsection{本章小结}
5.3 从 GMM-HMM 到 DNN Acoustic Model 这一部分把 Acoustic Model Overview 中相邻的教学单元串成了一条连续叙事：既交代了这一阶段要解决的问题，也交代了方法、推理或工程权衡为什么要这样组织。

\section{5.4 现代 Acoustic Model 架构}
根据课程页主题、公开视频字幕以及本讲的 coverage ledger，5.4 现代 Acoustic Model 架构 这一部分主要围绕 CNN、RNN、self-attention 作为 acoustic model architecture 的扩展路径。 展开。
写作上保持 coverage-first 原则：先把老师在这一阶段真正讲过的问题设定、模型部件、训练或推理逻辑讲清，再压缩成适合复习的结构。

\subsection{CNN、RNN、self-attention 作为 acoustic model architecture 的扩展路径}
从证据位置看，这一部分对应的课程证据区间是 2022-10-31 row。课程在这里强调的核心内容可以概括为：CNN、RNN、self-attention 作为 acoustic model architecture 的扩展路径。。
与单调模型不同，attention encoder-decoder 更强调全局条件化能力，但也因此引入了对齐自由度和部署层面的权衡。

记录系统中的附加说明指出：Taken from the official same-course 2022 row titled Advanced neural network architectures for acoustic model.

在本章结构中，这一单元被并入“5.4 现代 Acoustic Model 架构”，目的是把同一阶段的定义、机制和权衡集中讲清。

\subsection{本章小结}
5.4 现代 Acoustic Model 架构 这一部分把 Acoustic Model Overview 中相邻的教学单元串成了一条连续叙事：既交代了这一阶段要解决的问题，也交代了方法、推理或工程权衡为什么要这样组织。

\section{5.5 输出单元、词典与建模接口}
根据课程页主题、公开视频字幕以及本讲的 coverage ledger，5.5 输出单元、词典与建模接口 这一部分主要围绕 output units、pronunciation lexicon 与 acoustic model interface 之间的耦合关系。 展开。
写作上保持 coverage-first 原则：先把老师在这一阶段真正讲过的问题设定、模型部件、训练或推理逻辑讲清，再压缩成适合复习的结构。

\subsection{output units、pronunciation lexicon 与 acoustic model interface 之间的耦合关系}
从证据位置看，这一部分对应的课程证据区间是 2022 rows。课程在这里强调的核心内容可以概括为：output units、pronunciation lexicon 与 acoustic model interface 之间的耦合关系。。
从课程组织上看，这一部分承担的是把局部概念嵌回整条识别流水线中的作用，避免读者把公式或模块孤立地记忆。

记录系统中的附加说明指出：Reconstructed from official same-course system overview materials because the 2023 lecture lacks public artifacts.

在本章结构中，这一单元被并入“5.5 输出单元、词典与建模接口”，目的是把同一阶段的定义、机制和权衡集中讲清。

\subsection{本章小结}
5.5 输出单元、词典与建模接口 这一部分把 Acoustic Model Overview 中相邻的教学单元串成了一条连续叙事：既交代了这一阶段要解决的问题，也交代了方法、推理或工程权衡为什么要这样组织。

\section{总结与延伸}
Acoustic Model Overview 这一讲在整门《Speech Recognition and Understanding》课程中的位置，是把 Acoustic Model Overview 讲成一个可复用的建模模块。
从教材化角度看，本讲最重要的不是记住单一术语，而是理解它如何嵌入完整的 automatic speech recognition pipeline：输入语音如何被表示、模型如何被训练、解码如何得到最终转写，以及这些设计分别带来什么假设和权衡。
由于公开源以课程页和公开视频字幕为主、缺少官方 slide PDF，本章在正文里尽量把术语、方法和工程语境解释完整，并把缺失项留在 omission 和 manifest 里，而不是凭空补写。
\end{document}
